
\documentclass[12pt,onecolumn]{article}

% Encoding and typography
\usepackage[utf8]{inputenc}
\usepackage[T1]{fontenc}
\usepackage{lmodern}
\usepackage{microtype}

% Page layout
\usepackage[margin=1in]{geometry}
\usepackage{setspace}
\setstretch{1.08}

% Mathematics
\usepackage{amsmath,amssymb}

% Links and author identification
\usepackage[colorlinks=true,
            linkcolor=blue,
            citecolor=blue,
            urlcolor=blue]{hyperref}
\usepackage{orcidlink}

% Page headers
\usepackage{fancyhdr}
\pagestyle{fancy}
\fancyhf{}
\setlength{\headheight}{14pt}
\lhead{Space Was Never a Completed Container}
\rhead{John C. W. McKinley}
\cfoot{\thepage}

% Unnumbered title-page footnote
\newcommand\blfootnote[1]{%
  \begingroup
  \renewcommand\thefootnote{}\footnote{#1}%
  \addtocounter{footnote}{-1}%
  \endgroup
}

\title{\textbf{Space Was Never a Completed Container}\\[0.4em]
\large\textit{Quantum Eligibility, Relativistic Registration,\\
and the Ordered Appearance of Extension}}

\author{John C. W. McKinley\,\orcidlink{0009-0005-7097-5035}}

\date{October 2026}

\begin{document}

\maketitle

\blfootnote{\scriptsize Draft version. DOI placeholder to be supplied upon publication.}

\begin{abstract}

This paper advances an interpretation of physical reality in which space is not an independently completed container and the objects ordinarily imagined to occupy it are not permanently finalized classical substances. Nor is there a separate physical arrangement sitting behind observable events. What we call a physical configuration is the set of physical relations and conditions expressed through lawful interactions. Physical actuality consists of those interactions and the registered changes in physical relationships, while the governing laws specify the admissible relations and outcomes. The familiar picture of completed objects occupying completed space consequently introduces an ontology that is not established by physical measurement. Objects are identified through the stable, lawfully constrained relationships governing their interactions, and spatial extension is the geometric expression of those relationships rather than an additional physical entity within which they occur.

The argument follows from the clarified architecture of the Timeless Light Model, which distinguishes timeless lawful structure from the physical outcomes governed by that structure. Quantum mechanics describes lawful eligibility, including the amplitudes, probabilities, conservation principles, and interaction relations governing admissible outcomes. Special and general relativity govern relativistic registration, supplying the coherent temporal, causal, and geometric organization of registered physical relationships. These are complementary architectural roles rather than successive stages through which mathematical possibility is transformed into physical substance. Their combination permits physical reality to be understood through lawful interactions and relational conditions without introducing a permanently completed classical foundation beneath them.

The interpretation extends to ordinary macroscopic objects, observers, measuring instruments, motion, persistence, and distance. A rock is objectively real because the lawful organization identified as the rock continues to constrain physical interactions in mutually consistent ways. Its appearance, mechanical resistance, weight, and chemical properties are different expressions of that organization, not evidence requiring an additional completed rock-substance behind the interactions. The observer and the measuring apparatus belong to the same physical order and possess no privileged exemption from quantum governance. Similarly, the measurement of distance establishes an objective geometric relationship among physical systems without independently establishing a spatial material occupying the interval.

Two interpretive no-go arguments follow. The success of classical measurement does not license the inference to primitive classical substance, and the success of geometric measurement does not license the inference to an independently existing spatial container. These restrictions do not challenge the predictive content of classical mechanics, quantum mechanics, or relativity. They concern the additional ontology conventionally attached to successful descriptions. The positive interpretation identifies physical actuality with lawful interactions and registered outcomes, configurations with the relational conditions expressed through those interactions, and spatial extension with their relativistically organized geometry. The universe is therefore not fundamentally a collection of completed things occupying completed space or passing together through a universally completed present. It is a coherent lawful physical order expressed through actual interactions and their changing relationships. The present was never universal, and space was never a completed container.

\end{abstract}

\noindent\textbf{Keywords:} Timeless Light Model; lawful structure; lawful eligibility; relativistic registration; physical actuality; relational configuration; spatial ontology; macroscopic persistence; quantum measurement; relational geometry; interpretive no-go.

\clearpage
\tableofcontents
\clearpage

\section{Introduction: What Is Sitting There?}

Consider a billiard ball resting upon a billiard table. The ball has a measurable diameter, a recognizable surface, a definite mass, and a position relative to the table and its surroundings. If the ball is struck, its motion follows predictable relationships involving the applied force, the surface of the table, and the properties of the ball. Classical mechanics describes this familiar situation with extraordinary success, and nothing in the present paper calls that success into question. The question concerns the physical picture ordinarily attached to the description.

When imagining the billiard ball at rest, we tend to picture a completed physical object occupying a definite portion of independently existing space. The table is another completed object, and the room provides the larger spatial container within which both are situated. The scene appears to consist of objects that already exist in their finished physical forms, while the laws of physics determine what happens when they interact. This picture is sufficiently familiar that its ontological assumptions are rarely distinguished from the measurements that support the classical description.

Yet the fact that the ball can be measured, located, and predicted does not establish everything contained in that mental picture. It establishes that physical interactions involving the ball exhibit stable and coherent relationships. The properties attributed to the ball are identified through such relationships, and the predictions succeed because subsequent physical interactions continue to obey the governing laws. Whether an additional, independently completed object exists beneath that lawful physical organization is a different question.

One response is to replace the ordinary description of a ball with the more technical description of a physical configuration. The ball possesses internal structure, its constituents occupy relative positions, and its interactions depend upon the organization of those constituents. This is a useful advance in physical description, but it does not resolve the ontological question if the configuration is itself imagined as a completed arrangement of independently existing things. The original billiard ball has then merely been replaced by a more elaborate object, while the supposed physical foundation remains unchanged.

The question must therefore be carried beyond the language of configurations. What does the configuration actually identify? What is physically present when a rock rests on the ground, a wall obstructs passage, or a detector registers an interaction? In particular, what is meant by saying that a physical arrangement exists independently of the interactions through which its properties are expressed?

The present paper proposes that physical actuality consists of lawful interactions and registered changes in physical relationships. The governing laws determine which relationships and outcomes are admissible, while the conditions expressed through actual interactions constrain what happens in subsequent interactions. The stable organization of these relationships permits the identification of objects and the prediction of their behavior. A physical configuration is the description of that organization, rather than an additional physical substance concealed beneath it.

The same reasoning applies to space. Measurements establish distances, directions, positions, boundaries, and other geometric relationships, and the governing laws permit those measurements to be predicted with remarkable precision. These achievements establish the objective significance of spatial geometry. They do not, by themselves, establish a separate entity called space that exists as a completed container around or between physical systems.

This distinction becomes clearer when viewed through the clarified architecture of the Timeless Light Model, developed in \textit{Physics Has Always Been Unified: General Timelessness, Null Proper Time, and the Clarified Architecture} \cite{Unified2026}. That architecture distinguishes the lawful eligibility described by quantum mechanics from the relativistic registration of physical outcomes. Neither description requires a completed classical background standing outside the physical relationships being described.

The present investigation extends that architecture to the ordinary physical world. It examines the reality of objects and measuring instruments, the meaning of physical actuality, the status of configurations, and the geometric relationships conventionally described as space. The argument does not deny the billiard ball, its motion, or the distance between it and the edge of the table. It asks whether these objectively successful descriptions require the further assumption that permanently completed classical things occupy an independently completed spatial stage.

The central claim is that they do not. Physical reality is expressed through lawful interactions and the relational conditions governing those interactions. Objects, persistence, motion, and spatial extension are identifiable features of that physical order, while the familiar picture of completed objects sitting inside completed space adds an ontological structure that the measurements themselves do not supply.

\section{The Clarified Architecture of Physical Reality}

The architecture established in \textit{Physics Has Always Been Unified} begins by distinguishing the governing laws of physical reality from the physical events and relationships governed by those laws \cite{Unified2026}. This distinction is important because physical descriptions frequently employ the language of time, location, motion, and change without separating the behavior of physical systems from the status of the laws used to describe that behavior.

A governing law specifies which physical relationships and outcomes are admissible. Conservation principles, symmetries, coupling relations, amplitudes, probabilities, and the mathematical conditions governing interactions belong to this lawful structure. The law itself is not another physical event occurring alongside the events it governs. It does not wait for an interaction, travel from one system to another, or accumulate an elapsed time of its own. A mathematical law may contain temporal and spatial variables because the physical relationships it describes possess temporal and spatial properties, but the presence of those variables does not make the governing law an object located in space or an event progressing through time.

This distinction is the basis of general timelessness. As developed in the earlier McKinley papers, timelessness refers to the status of lawful structure rather than the denial of temporal relationships among physical events \cite{GeneralTimelessness2026,TimeNotGlobal2026}. A massive physical system accumulates proper time along its timelike worldline, and the temporal ordering of physical interactions retains its ordinary relativistic significance. The governing laws, however, do not themselves acquire a universal age or a position within the temporal relationships they organize.

An analogous distinction applies to spatial extension. The laws governing spatial relationships are not themselves spatially extended objects. A mathematical rule relating the measured separation of two physical systems to their interactions does not occupy the interval between them. Similarly, the equations describing gravitational geometry do not require the governing mathematical relationships to reside within a separate spatial environment. Their physical significance lies in the geometric and causal structure of the outcomes they govern.

Within this architecture, quantum mechanics describes lawful eligibility. Its mathematical formalism specifies the admissible outcomes of physical systems and the amplitude and probability relationships governing those outcomes. The wave function represents this eligibility structure, and the success of its predictions does not independently require that it be interpreted as an additional material object occupying the space ordinarily assigned to the physical system.

Special and general relativity provide the complementary account of relativistic registration. Physical outcomes possess causal, temporal, and geometric relationships that are organized according to the governing relativistic laws. The invariant propagation bound, proper-time accumulation, spacetime intervals, and the observer-consistent structure of registered physical events belong to this description.

The two roles are architectural rather than chronological. Lawful eligibility is not an earlier physical stage in which potential outcomes wait to become real, and relativistic registration is not a later manufacturing process through which mathematical possibilities are converted into material objects. Quantum mechanics and relativity describe different aspects of the same lawful physical order. Quantum mechanics supplies the eligibility structure, while relativity governs the spacetime organization of the registered physical relationships.

A physical registration is an actual interaction or outcome within that order. It does not require a conscious observer to witness the event. An atom interacting with another physical system, a detector undergoing a measurable change, a material surface responding to applied force, and a recording instrument preserving information about an interaction all participate in physical registration. The registering apparatus is itself a physical system governed by the same lawful structure as the system whose behavior it records.

This last observation has substantial consequences for the ontology of ordinary objects. If the detector, ruler, clock, laboratory, and observer are themselves quantum-governed physical organizations, they cannot serve as an independent classical foundation standing outside the relationships under examination. Their own stability and measurable properties require explanation within the same architecture.

The physical world is therefore not understood as quantum-governed activity occurring upon a completed classical stage. The same lawful structure governs the interactions of microscopic systems, macroscopic objects, measuring instruments, and observers. Their differing physical behavior is expressed through the organization and stability of their relationships, not through membership in fundamentally separate categories of physical existence.

\section{What Does Actual Mean?}

The distinction between lawful eligibility and registered actuality raises a question that is easily obscured by ordinary language. When we say that something is actual, what physical condition are we identifying?

Suppose a detector is prepared to register an interaction. Quantum mechanics describes the admissible outcomes under the applicable physical conditions and supplies the mathematical relationships governing their probabilities. These predictions have physical significance because they accurately constrain what the detector will register, but the mathematical prediction is not itself the occurrence of the predicted outcome. The actual detector response is a physical interaction whose registered consequences become part of the conditions governing subsequent interactions.

If the detector produces an electronic signal, that signal participates in further physical activity. It may be amplified, transmitted, stored, or used to change the state of another apparatus. A stored record may later be read by another physical system, and the record-bearing apparatus itself undergoes physical changes governed by the same laws. The actuality of the original interaction is expressed through the occurrence and consequences of these lawful changes, rather than through the arrival of an additional substance beneath them.

The distinction is between admissibility and occurrence. The governing laws establish the possible relationships and outcomes, while an actual physical interaction establishes the outcome that has occurred. The relational conditions expressed through that outcome then constrain the eligibility and organization of subsequent physical interactions.

This account also clarifies the meaning of objective reality. A rock is called real because its physical properties and interactions remain mutually consistent. Different instruments register compatible information concerning its location, boundaries, mass, composition, and mechanical behavior. The results do not depend upon an observer merely expecting to encounter a rock, since the laws governing the interactions are not controlled by the observer's beliefs. An incorrect expectation produces no authority over the physical outcome.

The continuing reliability of these interactions is what supports the identification of a persistent object. A rock does not become real because someone predicts that a subsequent measurement will again identify it. The prediction succeeds because the lawful physical relationships identified with the rock continue to constrain subsequent interactions in a manner consistent with its recognizable organization.

The difference between a physical fact and knowledge of that fact therefore remains important. An interaction occurs as part of physical actuality, whereas knowledge of the interaction consists of records, measurements, and descriptions that themselves have physical implementations. A photographic record, a detector reading, and a remembered observation are not the original interaction. They are later physical organizations whose relationships preserve information concerning it.

Recognizing this difference prevents the present account from reducing reality to perception or expectation. Physical outcomes are objectively constrained by the laws governing the interacting systems, and those constraints remain operative whether or not an observer understands or predicts them correctly. What changes in the proposed ontology is not the objectivity of physical behavior, but the assumption that such objectivity requires an independently completed classical substance beneath the interactions.

Physical actuality belongs to lawful interactions and registered outcomes. Stable objects are identifiable organizations of the relationships expressed through those outcomes, and descriptions of physical configurations organize our understanding of the conditions governing further interactions. No additional completed physical thing must be placed behind these relationships in order for their consequences to be real.

\section{A Configuration Is Not Another Physical Thing}

The term ``physical configuration'' is indispensable in ordinary physical description. A molecular configuration identifies the organization of a molecular system, a laboratory configuration identifies the relevant relationships among the apparatus and its surroundings, and a gravitational configuration describes the conditions governing the interactions of massive systems. Such descriptions are effective because the behavior of a physical system depends upon the relational conditions under which its interactions occur.

The ontological difficulty begins when the configuration is treated as though it were an additional physical object. In the case of a billiard ball, the description of its internal organization involves relationships among its constituents, the physical properties resulting from those relationships, and the conditions governing its interactions with the table and surrounding systems. To describe these features as a configuration is to organize the relevant physical information. It does not establish an independent configuration-substance existing behind the physical relationships being described.

Heating the ball changes its material organization and the conditions governing its subsequent interactions. A collision changes the momentum relationships involving the ball and the systems with which it interacts. Damage to its surface alters the interactions through which its shape, texture, and mechanical properties are registered. These changes are lawfully governed, and their outcomes contribute to the relational conditions relevant to future events.

It is unnecessary to imagine that a hidden, completed configuration persists beneath these changes as an additional constituent of the ball. The physical interactions and their relational consequences provide the relevant physical content, while the configuration describes that content in a form useful for analysis and prediction.

The same distinction becomes especially important when discussing the configuration of the universe. Ordinary language encourages us to imagine the universe as a completed arrangement of physical contents existing simultaneously at one universal moment. The phrase ``current configuration of the universe'' can inadvertently preserve this picture even after the notion of independent spatial substance has been rejected.

Relativity supplies no unique universe-wide present. Spatially separated events do not belong to one physically preferred universal simultaneity relation, and the information available to a local physical registration is restricted by relativistic causal structure. The earlier McKinley paper \textit{The Present Was Never Universal} develops the consequences of this local-registration interpretation \cite{Present2026}.

A mathematical description may assign physical variables across a selected spacelike hypersurface, and such descriptions remain useful in relativity and cosmology. Their usefulness does not establish that the chosen hypersurface corresponds to one independently completed universal present. The mathematical configuration organizes relationships under an adopted geometric description; it does not become another physical object standing behind the registered interactions.

The physical conditions relevant to an actual interaction are established within the lawful order governing that interaction. Other local registrations possess different causal relationships and different information availability, while the relativistic structure preserves the consistency of their descriptions. Consequently, neither physical actuality nor the meaning of configuration requires the entire universe to occupy one privileged completed state.

Persistence is understood in a similar manner. A stable physical organization remains identifiable because the governing relationships continue to support recognizable outcomes even as particular interactions change the system. The higher-order physical pattern persists through lawful continuity rather than through the permanent completion of each constituent in an unchanging classical arrangement.

A configuration is therefore a relational description of the conditions expressed through physical actuality. It does not identify a separate physical thing that undergoes change while remaining concealed behind physical events. The changes in lawful physical relationships are themselves the relevant physical events, and the language of configuration organizes the conditions under which those events occur and through which subsequent behavior is predicted.

\section{The Rock: Reality as Consistent Lawful Interaction}

A rock resting upon the ground provides a useful example because its existence seems entirely independent of the observer. We see its shape, color, and texture, encounter its mechanical resistance when touching it, register its weight when lifting it, and obtain further information through chemical analysis or other measurements. These interactions support the ordinary statement that a rock exists at a particular location, and the statement remains appropriate within the present interpretation.

The question is what makes that statement physically true. Under the clarified architecture, the rock is understood as a stable organization of lawful physical relationships. Its material properties, boundaries, and interactions are constrained by those relationships, while the continuing consistency of their outcomes preserves the recognizable physical pattern identified as the rock.

Visual observation involves electromagnetic interactions among the relevant physical systems, followed by registered changes within the observer's visual apparatus. Mechanical contact involves relationships that produce resistance and other measurable effects. A scale registers the behavior associated with the rock's weight, while a chemical instrument records information concerning its composition. These are not unrelated appearances assembled by an observer into an otherwise imaginary object. They are different physical interactions whose outcomes remain mutually consistent because they are governed by the same organized physical conditions.

Ordinary language gathers this consistency under the name ``rock.'' The name identifies the stable physical organization that supports the interactions, rather than supplying an additional substance beneath them. In this interpretation, the rock is not something distinct from its lawful physical organization that must separately explain why the organization produces rock-consistent results.

This account does not make the existence of the rock dependent upon human attention. Rain interacts with its surface, sunlight produces thermal and electromagnetic effects, and the surrounding ground participates in mechanical and gravitational relationships with it. Other physical systems encounter the rock's resistance or interact with its material structure whether or not a person is present to observe the event. The governing relationships remain objective because their physical consequences are not determined by human recognition.

The phrase ``the rock is there'' consequently expresses more than a present visual impression. It identifies a lawful physical organization that continues to constrain interactions and whose recognizable properties persist across different forms of measurement. An observer approaching the rock encounters the physical consequences of those relationships, rather than creating the rock by expecting an encounter.

The familiar image of a completed object simply sitting between observations nevertheless introduces an additional ontological assumption. It imagines that the rock possesses a permanently finalized classical existence beneath the physical relationships through which its properties are expressed. The fact that subsequent interactions continue to identify a rock establishes the reliability and persistence of its lawful organization, but it does not separately establish that additional substance.

The same reasoning applies to a wall. A wall blocks passage because the physical relationships governing its structure and its interaction with other systems produce mechanical resistance. Its ability to obstruct movement is neither subjective nor merely linguistic. The physical organization identified as the wall constrains the outcomes available when another physical system encounters it.

Ordinary material objects therefore remain objectively real in the present account. What changes is the explanatory interpretation of their reality. A rock or wall is not treated as an independently completed classical substance whose properties are subsequently expressed through interactions. Its identifiable existence consists in the stable, lawfully governed physical organization that those interactions reveal and through which its continued behavior is constrained.

\section{The Observer Is Another Lawful Physical Organization}

The rock example becomes more instructive when the observer is included in the physical description. An eye is not a detached classical spectator examining a finished object from outside the lawful physical order. It is itself a physical system composed of interacting constituents whose behavior is governed by electromagnetic, chemical, quantum, and other physical relationships.

When a person sees a rock, the optical structures and photoreceptors of the visual system participate in physical interactions associated with the rock and its surroundings. Those interactions produce changes that contribute to neural processing, recognition, memory, and subsequent behavior. The visual system is therefore not simply receiving an image of an independently completed object. It is participating in lawful interactions with another organized physical system.

The resulting registrations support the identification of the rock because the physical relationships involved remain mutually consistent. The rock contributes one stable lawful organization, the visual apparatus another, and their interactions produce physical changes whose outcomes preserve information about their relationship. Neither system occupies a privileged ontological position outside the laws governing both.

This distinction is not restricted to conscious observation. A photographic plate, detector, spectrometer, interferometer, or atomic clock participates in physical interactions without requiring conscious awareness to produce a registered response. The reliability of these instruments depends upon the lawful organization of their components and the stability of the relationships through which they interact with other systems.

For experimental purposes, it is useful to assign a measuring instrument a definite classical configuration. Its position, orientation, calibration, and recorded output can then be incorporated into the mathematical treatment of an experiment. That practice does not establish that the instrument has escaped quantum governance or acquired an independently completed classical substance. Its apparent classical stability is itself a macroscopic expression of lawful physical organization.

The distinction matters because descriptions of quantum measurement frequently contrast a quantum system with an apparatus treated as possessing definite classical properties. This division is useful for calculation, but it should not be converted without further justification into a division between two fundamentally different kinds of physical existence.

The measuring instrument is another physical participant whose interactions, records, and changes belong to the same lawful order as the system being measured. A detector registering a photon undergoes physical changes that participate in further interactions with its electronics, its recording apparatus, and any system that subsequently reads its output. Each part of the process remains subject to the governing physical laws.

The resulting chain does not require a privileged boundary at which quantum-governed relationships suddenly become independent classical substance. Instead, lawful physical organizations interact with one another, and their registered outcomes establish further conditions for subsequent interactions.

Seeing, measuring, and recording are consequently different expressions of physical participation within one coherent lawful structure. The ordinary distinction between an observer and an observed object remains useful when specifying the roles of systems in an experiment, but it does not imply that the observer is an intrinsically classical spectator or that the observed system must be an independently completed thing.

\section{No-Go on Primitive Classical Substance}

The preceding account supports an interpretive restriction upon the inference from successful classical description to primitive classical substance. The restriction follows from distinguishing what physical measurements establish from the additional ontology conventionally imagined beneath their results.

Consider repeated measurements of a billiard ball. Its mass, dimensions, position, and motion are identified with high reliability, and its behavior follows predictions derived from the governing physical laws. These measurements establish the existence of a stable and objectively constrained physical organization. They do not separately establish an additional completed classical substance whose existence is independent of the relationships producing the measurements.

The distinction becomes especially clear when the measuring instruments themselves are included in the account. Every instrument used to measure the billiard ball is another physical system whose properties and responses depend upon lawful interactions. The ruler has a stable material organization, the scale registers mechanical relationships, and the detector produces records through changes in its own physical state. The combined evidence establishes mutually consistent physical registrations, all governed by the same lawful structure.

To infer from this consistency that an additional primitive classical substance must exist beneath the relationships requires a further ontological step. The successful measurements do not independently supply that step. The reality of the ball as a stable physical organization is already supported by its lawful interactions, and the proposed classical substance adds no separately established physical content.

This is the first interpretive no-go. Its conclusion concerns what may legitimately be inferred from the measurements, rather than a mathematical demonstration that every conceivable metaphysical substrate is impossible. The absence of a required inference is sufficient for the present architectural purpose, because primitive classical substance need not be introduced as a fundamental constituent of the physical account.

The reasoning is related to the earlier McKinley paper \textit{Wavefunction Prediction Does Not License a Photon Path} \cite{WavefunctionNoGo2026}. In that treatment, successful wavefunction prediction does not establish an independently existing intermediate photon trajectory. The photon argument also invokes the specific relativistic properties of null propagation, including the absence of a photon rest frame and the vanishing of proper-time accumulation along a null path.

Those photon-specific properties are not transferred to macroscopic matter. A billiard ball or rock participates in timelike physical relationships, and the corresponding descriptions of proper-time accumulation, trajectories, momentum, and motion remain valid. What the two arguments share is the distinction between successful physical representation and the additional ontology inferred from it.

For the billiard ball, the registered physical evidence establishes a stable macroscopic organization whose interactions obey predictable relationships. Classical mechanics describes that organization accurately, but the accuracy of the description does not independently establish a permanently finalized classical thing beneath it.

The positive account therefore retains the objectively real billiard ball while declining to make primitive classical substance a necessary part of its physical ontology. The ball remains identifiable because its lawful relationships continue to constrain its interactions, not because a second, completed ball must be imagined behind the physical activity through which its properties are expressed.

\section{Space Is Not Something Sitting Between Objects}

The familiar conception of space contains an assumption closely related to the conception of primitive classical substance. When two rocks are separated by a measurable distance, we ordinarily picture them occupying different locations within an already existing spatial container. The interval between them is imagined as a portion of that container, which remains independently present whether or not any measurement occurs.

This picture is useful for ordinary geometric reasoning, but its usefulness does not establish the existence of the container as an additional physical entity. What physical measurement establishes is a relationship among physical systems. The geometric character of that relationship is objective, and the result is constrained by the governing laws, but the measurement does not separately identify a spatial material occupying the interval.

Suppose a ruler is placed between the rocks. The ruler has its own stable physical organization, and the interactions involved in positioning and reading it establish information concerning the separation of the rocks. The measurement succeeds because the relevant physical relationships possess consistent geometric properties. The ruler, the rocks, and the measuring apparatus participate in interactions that support a determinate distance.

Radar ranging uses a different physical procedure to establish a related geometric quantity. Clock readings and signal interactions are combined according to the governing laws to calculate the separation of physical systems. Interferometry likewise supplies geometric information through phase relationships and registered detection outcomes. These methods differ in their physical implementations, but they share the feature that the resulting spatial quantities are established through lawful relationships among physical systems.

The success of these measurements establishes that spatial relationships are real and geometrically constrained. It does not require a further substance called space to be located between the systems. The measured distance is already part of the lawful physical description, and the additional spatial container contributes no separately identified measurement.

The distinction is not merely that empty space has never been photographed as a material object. That observation alone would be insufficient, since many physically meaningful quantities are established indirectly. The argument instead concerns the explanatory content of the physical architecture. Spatial relationships are assigned to the lawful geometric organization of physical interactions and outcomes, so a completed container performs no additional required role in explaining what is measured.

A coordinate identifies a position within a mathematical representation of physical relationships. It does not establish an independently completed physical object corresponding to that coordinate. Similarly, a measured interval establishes a geometric relationship rather than an additional substance filling the interval. The prediction of an interaction across a distance confirms the governing laws relating the physical systems involved, without requiring a separate spatial stage upon which the interaction is performed.

Under this interpretation, space is the geometric expression of lawful physical relationships. Objects themselves are identified through stable patterns of lawful organization, while their positions and separations describe geometric aspects of the relationships governing their interactions.

This does not eliminate distance or weaken its physical significance. A meter remains an objective unit of measured extension, and the laws governing its application retain their predictive authority. What is rejected is the additional identification of that meter with a portion of independently completed spatial substance.

The two rocks remain separated by a measurable geometric relationship. The measurement is real, the relationship is physically consequential, and its predictions are governed by established physical laws. What is absent from the proposed ontology is a further spatial entity sitting between the rocks as a completed container.

\section{Spatial Geometry as the Lawful Description of Relationships}

The rejection of an independently completed spatial container does not entail the rejection of geometric structure. Special and general relativity retain their full mathematical and predictive significance because the present interpretation concerns what the geometric descriptions establish about physical reality, rather than proposing alterations to the equations themselves.

In special relativity, the invariant structure of spacetime intervals permits different observers to employ different coordinate descriptions while preserving the relevant physical relationships. Proper-time accumulation, causal ordering, and the invariant information-propagation bound are not discarded when coordinates are recognized as descriptions rather than independently existing physical substances. They remain features of the lawful geometric organization of physical interactions.

General relativity supplies a more extensive account in which the metric determines the local spacetime interval,

\begin{equation}
ds^2 = g_{\mu\nu}dx^\mu dx^\nu.
\end{equation}

The metric also participates in the mathematical description of gravitational relationships and curvature, with the Einstein field equations expressing the relationship between spacetime geometry and stress-energy,

\begin{equation}
G_{\mu\nu}+\Lambda g_{\mu\nu}
=
\frac{8\pi G}{c^4}T_{\mu\nu}.
\end{equation}

Nothing in the present argument changes these equations or the experimental predictions derived from them. The issue is whether their success establishes an additional physical substance corresponding to an independently completed spatial container.

The metric provides the geometric relationships governing intervals, causal structure, proper-time accumulation, and gravitational behavior. Its physical significance is confirmed through the accuracy with which those relationships describe and predict interactions. The fact that geometric quantities have measurable consequences does not, however, independently establish a material container underlying the mathematical description.

This distinction is particularly important in general relativity, where geometry is not merely a passive background of the Newtonian kind. Curvature has measurable effects, gravitational waves produce registered responses in suitable detectors, and solutions of the Einstein field equations possess meaningful geometric structure even where the stress-energy tensor vanishes. The present interpretation preserves these features by treating geometry as physically consequential lawful structure rather than identifying it with an additional completed spatial substance.

A gravitational-wave detector, for example, registers changes in physical relationships among its components according to the relevant relativistic geometry. The observed response confirms the physical consequences of that geometry. It does not require the detection of a second material constituent called space beneath the relationships whose changes are measured.

The distinction is equally evident in the operational determination of distance. Consider an idealized radar-ranging measurement in which an observer sends a signal toward a reflector and receives the returning signal. Under the usual approximately flat and stationary conditions, the radar distance is expressed as

\begin{equation}
D_{\mathrm{radar}}
=
\frac{c}{2}
\left(
\tau_{\mathrm{return}}
-
\tau_{\mathrm{send}}
\right).
\end{equation}

The proper-time readings belong to the observer's clock, and \(c\) is the invariant relativistic propagation bound. The distance is established through lawful physical relationships involving the clock, emitter, signal, reflector, receiver, and the registered interactions connecting the measurement procedure. More general geometries require the corresponding metric treatment, but the conceptual point remains the same: the physical measurement determines a geometric relationship rather than detecting an independently existing spatial substance.

Directions and boundaries are understood within the same framework. An optical apparatus registers directional information through the lawful relationships governing its interactions. A material boundary is identified through the physical properties and interactions that distinguish the behavior of a system from that of its surroundings. The boundary of a wall, for example, has practical and measurable significance because interactions with the wall produce the resistance and other outcomes predicted from its lawful material organization.

The reality of that boundary does not establish a second substance called space on either side of it. The boundary is part of the geometric and physical description of the relevant relationships, while the behavior associated with it remains governed by the physical laws.

Spatial extension is therefore retained as an objective feature of the physical description, but it is not assigned the role of an independent container. Extension belongs to the geometrically organized relationships whose registered consequences are governed by relativity. The geometry is physically significant because it constrains interactions, not because it must be imagined as an additional object within which those interactions occur.

\section{No-Go on Independent Spatial-Container Ontology}

The analysis of spatial measurement supports a second interpretive restriction, parallel to the restriction concerning primitive classical substance. The successful measurement and prediction of geometric quantities do not independently establish that an additional completed spatial container exists beneath the relationships being measured.

Physical determinations of distance, position, direction, curvature, and extension involve interactions among physical systems. Their outcomes are interpreted through mathematical descriptions whose consistency is governed by established physical laws. The success of these descriptions confirms the physical significance of the geometric relationships, but it does not independently identify a further material entity corresponding to space itself.

The inference from measured geometry to independently completed spatial substance therefore requires an ontological assumption beyond the physical result. A coordinate system organizes physical relationships without becoming a collection of completed physical locations waiting to be occupied. A metric describes lawful geometric structure without automatically becoming material substance, and an empty region described by a relativistic solution need not be interpreted as an object composed of spatial material.

Within the clarified architecture, the governing laws supply the geometric and causal structure of physical relationships, while registered interactions express those relationships in actual physical outcomes. The successful physical description already contains the measured geometric content. Introducing a separate container does not add a further independently established relationship or registered outcome.

This is the second interpretive no-go. As with the first, the restriction concerns the necessity of an ontological inference rather than a mathematical disproof of every conceivable substantival theory. The measurements establish the objective geometry, but they do not compel the additional conclusion that geometry must belong to an independently completed spatial substance.

The positive interpretation is accordingly more specific than a general skepticism about the existence of space. It retains spatial measurements, geometrical relationships, and the complete predictive structure of relativity while identifying space with the geometric expression of lawful physical relationships. The physical content ordinarily attributed to an independent container is already supplied by those relationships and their registered consequences.

The parallel with macroscopic objects is instructive. A rock is real as a stable lawful physical organization, rather than as an independently completed classical substance added beneath its interactions. A distance is real as an objective geometric relationship, rather than as a portion of an independently completed spatial material added beneath its measurement.

In both cases, the physical description remains intact while the additional ontology is withheld. The stable object and the measured separation are not discarded. What disappears is the supposed requirement for another independently completed physical constituent standing behind them.

\section{Motion, Persistence, and the Absence of a Completed Stage}

The ordinary classical picture combines completed objects and completed space with a further assumption concerning time. A billiard ball is imagined as a persisting object that occupies one position, moves to another, and continues to a third while a universal temporal stage advances. The resulting picture is useful because the familiar language of trajectories and motion accurately describes a wide range of physical behavior.

The clarified architecture retains those descriptions while interpreting their physical content through lawful changes in relationships. A macroscopic system possesses stable physical organization, and its interactions with other systems produce changes whose temporal and geometric ordering is governed by relativity. Its recognizable identity persists because the relevant physical relationships continue to support a coherent macroscopic description. The trajectory represents the ordered changes in those relationships, without requiring a permanently finalized object to traverse an independently completed spatial substance.

This approach is consistent with the earlier treatment of temporal structure in the Timeless Light Model. The null photon case demonstrates that proper-time accumulation is not a universal property of all physical propagation. The absence of a photon rest frame and the null character of its proper-time interval prevent the ordinary traveler narrative of a massive body from being transferred literally to photon propagation \cite{NullProperTime2025,WavefunctionNoGo2026}.

Massive physical systems, by contrast, possess timelike relationships and accumulate proper time along their worldlines. The present interpretation does not diminish that physical distinction. Instead, it recognizes that temporal properties belong to the physical relationships in which those properties occur, without requiring the governing laws themselves to possess a universal temporal variable.

Spatial structure is treated analogously. A physical relationship may possess measurable extension, direction, or curvature, and relativity governs the corresponding geometry. These properties do not require the governing lawful structure to occupy an independent spatial container. Temporal and geometric quantities remain physically meaningful where the relevant relationships supply temporal and geometric structure.

The meaning of persistence becomes clearer under this account. A rock can remain identifiable across many interactions because its lawful physical organization continues to preserve the relationships supporting its recognizable properties. Its continued existence does not require every microscopic constituent to remain permanently finalized in a classical state, nor does it require an additional completed rock to remain concealed beneath changing interactions.

A radioactive material illustrates this distinction. Its governing physical laws specify admissible decay processes and the probabilities associated with them. Individual decay events occur, and their registered consequences alter the relevant material conditions. The larger sample may nevertheless retain a recognizable macroscopic organization throughout many such events. Continuity of identity is supported by the stability of the relevant higher-order relationships, even while particular constituents undergo physical changes.

A wall provides a more familiar example. Its constituent physical systems participate in continuing lawful activity, while the macroscopic organization remains sufficiently stable for its location, shape, and resistance to be reliably predicted. The wall persists because its governing relationships continue to support wall-consistent interactions, not because its physical identity must be supplied by a permanently completed classical substance beneath those relationships.

Motion is understood in the same way. A moving object has changing geometric relationships with its surroundings, and those changes are temporally and causally ordered according to the governing laws. Classical mechanics represents this behavior through quantities such as position, velocity, acceleration, and trajectory. Its success follows from the stability and predictability of the underlying physical relationships.

The classical account remains appropriate because it provides an efficient description of the observed physical behavior. What the present interpretation declines to infer from that success is the existence of a completed classical object crossing a completed spatial stage. Motion is the lawful change in spatial relationships involving an identifiable physical organization, and persistence is the continuing lawful consistency of that organization across its interactions.

\section{The Universe Has No Completed Universal Configuration}

The preceding analysis has consequences for the conception of the universe as a whole. If objects are identified through stable lawful physical relationships, if spatial extension expresses the geometry of those relationships, and if temporal ordering is not governed by one universal present, then the universe should not be pictured as a single completed arrangement of objects occupying a completed spatial container at a universally shared moment.

Such a picture imports the classical notion of a universal snapshot into a physical architecture that does not supply one. Relativity permits different coordinate descriptions and different assignments of simultaneity to spatially separated events. It does not establish a unique physical present extending across the entire universe independently of those descriptions.

Local physical registrations possess their own causal relationships, and the physical information available at one location differs from the information available at another. These differences are not failures of universal coordination that must be corrected by introducing a hidden common present. They belong to the relativistic organization of physical events, under which information availability and causal relationships are constrained by the governing geometry.

The notion of a current configuration of the universe must therefore be distinguished from the mathematical methods used to describe physical systems. A relativistic model may assign fields, coordinates, or other variables across a selected spacelike hypersurface. Such representations are useful in general relativity, cosmology, and related calculations, but their mathematical validity does not establish that the selected hypersurface is the independently existing present of the universe.

The representation organizes relationships under a particular geometric description. It does not create an additional physical configuration standing behind the local registrations, nor does it require all physical systems to occupy one permanently finalized classical arrangement.

This distinction also applies to the reconstruction of physical history. An observer may use presently available records to infer earlier events and develop an account of how a system acquired its observed properties. A geological formation, photograph, detector memory, or clock record contains physical information that supports such inference. The record itself, however, is a physical organization participating in presently available interactions, rather than the earlier event reproduced as an independently existing physical object.

The success of historical reconstruction does not require the earlier events to remain physically completed within a universal temporal archive. Similarly, the success of a calculated spatial configuration does not require the described arrangement to exist as an additional completed object occupying the coordinates of the calculation.

The present argument does not deny the coherence of relativistic models, the value of reconstructed physical histories, or the usefulness of describing wide regions of the universe. It distinguishes those descriptions from the ontology claimed on their behalf.

The universe is governed by coherent lawful structure, while actual physical interactions and registered changes possess the causal and geometric relationships specified by that structure. The conditions relevant to further interactions arise within this lawful physical order, without requiring one universal configuration to stand behind them as a completed physical snapshot.

Replacing completed space with a completed universal configuration would merely preserve the earlier assumption under a different name. The clarified architecture instead treats configurations as descriptions of lawful relationships, and it treats physical actuality as the interactions and registered changes through which those relationships are expressed.

\section{Classical Mechanics as a Compression of Lawful Relationships}

The rejection of primitive classical substance and independently completed spatial containers does not impair the practical usefulness of classical mechanics. A billiard ball retains its measurable mass, position, and velocity, a table supports material objects, a rock resists mechanical contact, and planetary behavior remains describable through the applicable gravitational laws. These descriptions succeed because the underlying physical relationships produce stable macroscopic regularities.

Classical mechanics organizes those regularities through a comparatively small number of variables. The complicated interactions among the constituents of a billiard ball, the table, and the surrounding environment can be represented through the position and motion of the ball, the relevant contact forces, and other quantities appropriate to the calculation. The resulting description is an efficient compression of physical relationships whose higher-order behavior remains stable.

The term ``billiard ball'' identifies one recognizable physical organization, while ``table'' identifies another. Their separation is a geometric relationship, and their collision is an actual interaction governed by physical laws. Classical mechanics describes the system through those familiar categories without requiring the microscopic details of every contributing physical interaction to be included in the calculation.

The simplification is justified by the reliability of the macroscopic behavior. A wall can be treated as rigid for many purposes, even though its constituent physical systems participate in continuing activity. A ruler can be assigned a stable length under appropriate measurement conditions, and a rock can be treated as having a definite location even though its internal physical organization is not a permanently frozen classical arrangement.

These practices succeed because the relevant physical relationships retain sufficient stability for the intended description. Their success does not imply that macroscopic systems are fundamentally exempt from the laws governing their constituents.

A useful comparison appears in thermodynamics, where collective physical behavior is accurately described through macroscopic variables that do not have to be assigned to every microscopic constituent in the same manner. Temperature and pressure provide effective descriptions of physical systems whose microscopic behavior is more complicated than the macroscopic variables themselves. Their usefulness arises from the organized behavior of the system rather than from the requirement that the macroscopic description identify a separate primitive substance.

The present interpretation applies a comparable discipline to classical objects and spatial geometry. The concept of a persistent object is an effective description of stable lawful physical organization, while the image of a spatial container is an effective way of visualizing geometrically organized relationships. Both forms of description assist reasoning and calculation, but their success does not establish that the associated mental pictures correspond to independently completed ontological primitives.

A room illustrates the distinction. Ordinary thought treats the room as a completed volume containing furniture and other objects. This picture is useful for navigation, construction, and measurement, but the physical measurements establish relationships among material boundaries, instruments, and other physical systems. The geometry of those relationships provides the measured extension without requiring a further substance called room-space.

In the same manner, a billiard ball moving across a table is accurately represented as a recognizable object following a trajectory through space. The physical content consists of the lawful organization identified as the ball, its changing relationships with the surrounding systems, and the temporal and geometric ordering of those changes.

Classical mechanics remains the appropriate language for the stable macroscopic behavior it describes. What the clarified architecture rejects is not that language, but the assumption that each noun appearing in an effective physical description must correspond to an independently completed physical constituent.

The ball, table, wall, and planet are real because their lawful physical organizations produce objective and predictable interactions. Their spatial relationships and motions are equally real because the governing laws constrain their measurement and behavior. The classical description is successful precisely because it captures stable aspects of this lawful physical order.

\section{Relationship to Existing Physics}

The present argument introduces no new equations, particles, forces, or physical constants. It changes none of the mathematical predictions of quantum mechanics, special relativity, general relativity, or classical mechanics. Its contribution concerns the interpretation of what these theories establish about physical reality and, in particular, the ontological assumptions that accompany their successful descriptions.

The underlying question has substantial precedent in the development of modern physics. Einstein's treatment of general covariance and physically meaningful spacetime relationships established important distinctions between coordinate representations and invariant physical content \cite{Einstein1916}. Subsequent relational approaches to quantum mechanics have examined the significance of interactions among physical systems without granting a single observer a privileged position outside the physical description \cite{Rovelli1996}.

Research on decoherence has likewise explained important features of the emergence and stability of effectively classical behavior in systems governed by quantum mechanics \cite{Zurek2003}. These developments provide relevant context for understanding how macroscopic physical descriptions remain successful without requiring a separate domain of fundamentally classical objects.

The present paper does not claim that the general idea of relational physical description is new. Its contribution lies in extending the clarified Timeless Light Model architecture to the interpretation of physical actuality, macroscopic object identity, configurations, and spatial extension.

The earlier architecture assigns complementary explanatory roles to quantum eligibility and relativistic registration. Quantum mechanics describes the admissible outcomes and their mathematical relationships, while relativity supplies the coherent causal, temporal, and geometric structure of registered physical outcomes. The present paper asks what additional physical ontology is necessary when those registered interactions and their lawful relational conditions are recognized as the relevant physical content.

The resulting interpretation declines to introduce independently completed classical objects or a completed spatial container because neither is separately established by the measurements or required to supply the physical relationships described by the existing theories.

The two interpretive no-go arguments identify the corresponding restrictions upon inference. The success of classical measurements establishes stable and predictable macroscopic physical organization, but not the additional necessity of primitive classical substance. The success of spatial and geometric measurements establishes objective lawful relationships, but not the additional necessity of an independently completed spatial material.

The proposed ontology also avoids interpreting registration as a hidden mechanism by which possibilities are manufactured into substances. Lawful eligibility is not a preceding physical realm, and registration does not require a separate selector operating outside the governing physical structure. The distinction between the laws and their registered outcomes is explanatory rather than chronological.

Physical interactions occur within the lawful order, and their registered outcomes possess the relativistic relationships specified by that order. Classical objects and familiar spatial relationships are identifiable stable expressions of this activity, without requiring the completed classical world to exist as a separate foundational stage.

\section{Discussion: What Remains When the Completed World Is Removed?}

The principal challenge of the present interpretation is to explain physical reality without unintentionally reconstructing the classical substrate it seeks to remove. It is easy to reject an independently existing spatial container while retaining permanently completed classical objects. It is equally easy to abandon primitive classical objects while replacing them with a complete physical configuration that is imagined as a new kind of substance behind observable events.

Neither replacement answers the underlying question. If the configuration is treated as a completed arrangement of independently existing constituents, the familiar classical ontology survives under a more technical name. If lawful relationships are pictured as literal threads connecting the constituents of a physical network, those threads may become another imagined substance suspended within an independently completed spatial environment.

The relational language of the clarified architecture is not intended to introduce such substitutes. A relationship is not a miniature cable extending between completed objects, and a physical configuration is not an invisible mesh occupying space. These expressions describe lawful physical conditions and the organization of actual interactions rather than a further material layer beneath them.

A network image nevertheless provides a useful illustration if its representational character is maintained. One network may depict the physical organization identified as a rock, while another depicts the physical organization identified as an eye or detector. The relationships within each network represent the lawful organization of the relevant system, and the relationships between them represent the conditions governing their interactions.

The image conveys that neither the rock nor the observing apparatus must be treated as a detached primitive classical object. It also illustrates why observations remain mutually consistent: the outcomes of interactions are constrained by the governing physical relationships of the participating systems.

The network itself, however, is only a representation. The proposed physical ontology contains no literal threads suspended within a completed container and no additional web standing behind actual interactions. Its physical content consists of lawful relations, their registered outcomes, and the conditions governing subsequent activity.

This distinction helps clarify why the interpretation does not reduce physical reality to a private construction of consciousness. The lawful relationships do not depend upon an observer's expectations, and their consequences occur through interactions among physical systems whether or not anyone is attending to them. The rock remains objectively consequential because its organization continues to govern the outcomes of physical interactions. The observer participates in that same order rather than supplying the order through perception.

The meaning of physical existence is consequently tied to the lawful physical behavior of the system. When a wall resists passage, the physical interactions are real and the resistance follows from the wall's material organization. When a planet produces measurable gravitational consequences, those consequences express the governing lawful relationships identified with the planetary system. When a detector registers a physical event, its response becomes part of the conditions governing later interactions.

The independent classical substance ordinarily imagined beneath these phenomena is not what the measurements establish. The registered outcomes and their lawful relationships are the physical facts through which the stability of the world is known and described.

Spatial measurement presents the corresponding situation. The interval between two systems is real as a lawful geometric relationship, and the measurement of that interval has objective predictive significance. The independent container ordinarily imagined beneath the geometry is not the measured result. Its introduction adds a picture of physical reality beyond the relationships established by the measurement.

The same principle governs the treatment of time and configurations. A reconstructed history organizes physical information preserved in records, but the reconstruction does not require past events to remain as completed physical objects in an independently existing temporal archive. A calculated configuration organizes information about lawful physical conditions, but the calculation does not establish a separate configuration-substance. A coordinate representation organizes spatial relationships without requiring a completed physical lattice corresponding to every coordinate.

Descriptions serve to organize information about physical reality. Their usefulness and predictive success do not automatically establish additional physical constituents corresponding to the imagery associated with those descriptions.

What remains after the completed classical world is removed is therefore not an absence of physical reality. The governing laws continue to specify admissible interactions and outcomes. Actual interactions occur, their registered consequences establish relational conditions, and those conditions participate in the lawful organization of subsequent physical activity. Stable physical patterns remain identifiable as objects, while their changing relationships are described as motion, persistence, and spatial extension.

The resulting picture preserves the objective character of the ordinary world while declining to multiply its fundamental physical constituents. A rock is not less real because no additional rock-substance is placed beneath its lawful organization. A measured distance is not less real because no separate spatial material is placed beneath the interval. The familiar world remains available to measurement, prediction, and interaction because those measurements and interactions are themselves the lawful physical activity from which the familiar descriptions arise.

\section{Conclusion: Space Was Never a Completed Container}

The familiar picture of the universe presents completed objects arranged within completed space and changing through a common temporal progression. This picture is effective for ordinary reasoning and for the classical descriptions of stable macroscopic behavior, but its effectiveness does not establish that the pictured objects, container, and universal stage must exist as independent foundational constituents of physical reality.

The clarified architecture of the Timeless Light Model provides a different interpretation. Governing laws specify admissible physical relationships, quantum mechanics describes the eligibility of physical outcomes, and special and general relativity govern the causal, temporal, and geometric organization of registered interactions. Physical actuality is expressed through those lawful interactions and the changes in relational conditions that accompany them.

This architecture permits the reality of ordinary objects to be understood through the stable physical organization governing their interactions. A rock remains objectively identifiable because its material relationships continue to constrain subsequent outcomes in mutually consistent ways. A measuring instrument participates in the same lawful physical order and possesses no privileged classical exemption. The continued identity of the rock or instrument is expressed through the stability of its physical organization, even while particular interactions and conditions change.

The first interpretive no-go follows from the distinction between these established physical relationships and the additional classical substance conventionally imagined beneath them. Successful classical measurement confirms the existence of stable, predictable macroscopic physical behavior, but does not independently require a permanently finalized classical object underlying the interactions.

The second no-go extends the same reasoning to spatial geometry. Measurements of distance, direction, curvature, and extension establish objective geometric relationships whose behavior is governed by relativity. Their success does not independently establish a completed spatial container occupying the intervals represented by those measurements.

The resulting interpretation retains the physical content of the existing theories while withholding an unnecessary additional ontology. A configuration describes the lawful relational conditions expressed through physical interactions, rather than identifying another completed thing concealed behind them. Spatial extension describes the geometry of those relationships, rather than a separate substance within which they occur. Persistence is the continued lawful consistency of an identifiable physical organization, while motion is the lawfully ordered change in its geometric relationships with other systems.

These distinctions also prevent the completed spatial container from being replaced by a completed universal configuration. Relativistic registration does not establish a unique universe-wide present in which all physical systems occupy one privileged, independently finalized arrangement. Local interactions possess the causal and geometric relationships supplied by the governing relativistic structure, and descriptions of wider configurations organize those relationships without creating an additional universal physical snapshot.

The universe is therefore not fundamentally a collection of independently completed things placed within independently completed space. Its physical actuality consists of lawful interactions and registered changes, with the governing structure specifying the relationships and outcomes through which stable physical behavior is expressed. The classical vocabulary of objects, locations, motion, and persistence remains valid because it describes objective and predictable features of that lawful activity.

The ordinary world has not been eliminated by this interpretation. Its physical content has instead been identified with the interactions and relationships that measurement has always established. Nothing further must be imagined as a completed classical substance beneath a rock, a completed configuration beneath physical events, or an independently existing spatial material between objects.

The present was never universal. Space was never a completed container.

\clearpage

\begin{thebibliography}{99}

\bibitem{Unified2026}
J.~C.~W.~McKinley,
\textit{Physics Has Always Been Unified: General Timelessness, Null Proper Time, and the Clarified Architecture---Why Quantum Mechanics and Relativity Were Never in Competition}.
Zenodo, September 12, 2026.
\href{https://doi.org/10.5281/zenodo.20954931}
{doi:10.5281/zenodo.20954931}.

\bibitem{WavefunctionNoGo2026}
J.~C.~W.~McKinley,
\textit{Wavefunction Prediction Does Not License a Photon Path: A Short Interpretive No-Go}.
Zenodo, April 12, 2026.
\href{https://doi.org/10.5281/zenodo.19504771}
{doi:10.5281/zenodo.19504771}.

\bibitem{TimeNotGlobal2026}
J.~C.~W.~McKinley,
\textit{Time Is Not Global: Disambiguating the Timeless Light Model from Wheeler--DeWitt, Page--Wootters, and Relational Quantum Time}.
Draft manuscript, September 2026.
\url{https://johncwmckinley.com/}.

\bibitem{Present2026}
J.~C.~W.~McKinley,
\textit{The Present Was Never Universal: Local Registration and the Absence of a Universe-Wide Now}.
Draft manuscript, July 2026.

\bibitem{GeneralTimelessness2026}
J.~C.~W.~McKinley,
\textit{Timelessness Is the General Condition: Spacetime Is the Special Case}.
Zenodo, 2026.
\href{https://doi.org/10.5281/zenodo.19771924}
{doi:10.5281/zenodo.19771924}.

\bibitem{NullProperTime2025}
J.~C.~W.~McKinley,
\textit{Taking Null Proper Time Seriously: An Interpretive Clarification of Null Proper Time}.
Zenodo, 2025.
\href{https://doi.org/10.5281/zenodo.18004632}
{doi:10.5281/zenodo.18004632}.

\bibitem{Rovelli1996}
C.~Rovelli,
\textit{Relational Quantum Mechanics}.
International Journal of Theoretical Physics
\textbf{35}, 1637--1678, 1996.
\href{https://doi.org/10.1007/BF02302261}
{doi:10.1007/BF02302261}.

\bibitem{Zurek2003}
W.~H.~Zurek,
\textit{Decoherence, Einselection, and the Quantum Origins of the Classical}.
Reviews of Modern Physics
\textbf{75}, 715--775, 2003.
\href{https://doi.org/10.1103/RevModPhys.75.715}
{doi:10.1103/RevModPhys.75.715}.

\bibitem{Einstein1916}
A.~Einstein,
\textit{Die Grundlage der allgemeinen Relativit{\"a}tstheorie}.
Annalen der Physik
\textbf{49}, 769--822, 1916.

\end{thebibliography}

\end{document}